\documentclass[11pt]{article}
\usepackage[margin=1in]{geometry}
\usepackage{amsmath,amssymb,amsthm}
\usepackage{enumitem}

\newtheorem{theorem}{Theorem}[section]
\newtheorem{lemma}[theorem]{Lemma}
\newtheorem{proposition}[theorem]{Proposition}
\newtheorem{conjecture}[theorem]{Conjecture}
\newtheorem{remark}[theorem]{Remark}

\newcommand{\Om}{\Omega}
\newcommand{\R}{\mathbb R}
\newcommand{\dT}{\delta_T}
\newcommand{\A}{\mathcal A}
\newcommand{\X}{\mathcal X}
\newcommand{\Et}{E_{\hat t}}

\title{Pushing the Planar Graph-Reduction Step}
\author{}
\date{May 2026}

\begin{document}
\maketitle

\section{Executive summary}

This note tries to prove the remaining planar reduction-to-graph step.  I do
not get a complete proof.  I do get a sharper formulation that looks more
realistic than the previous \(O(\dT)\)-radialization target.

The key adjustment is to stop at a boundary layer of size
\[
   \lambda \simeq \sqrt{\dT(\Om)}
\]
rather than at size \(O(\dT(\Om))\).  This is still sharp for the final
asymmetry estimate, because the transfer back costs \(O(\lambda)\) in
Fraenkel asymmetry and hence \(O(\lambda^2)=O(\dT)\) after squaring.  The gain
is that the level-set Chebyshev step now gives
\[
   D_H(\hat t)+D_I(\hat t) \lesssim \sqrt{\dT(\Om)}
\]
instead of only a fixed small constant.  In dimension two this gives real
geometric information: small components have total area \(O(\dT)\), holes with
small perimeter have total area \(O(\dT)\), and connected outward fingers of
macroscopic length are excluded by perimeter unless they are exactly the thin
objects that the \(D_H\) weight detects.

The proof still has one missing arbitrary-geometry lemma: a
\emph{good-level convexification/radialization lemma}.  It must say that
convexifying, filling, or pruning only the non-graph pieces of the good level
changes the normalized torsion deficit by \(O(\dT)\).  The model estimates
support this, but I do not see a complete decomposition proof in the current
repo.

\section{The better stopping scale}

Let \(|\Om|=\pi\) and set \(\delta=\dT(\Om)\).  The volume-variable identity
from Plan 3 Step 1 gives, for \(s=m(t)\),
\[
   \int_0^\pi
      \frac{D_H(v(s))+D_I(v(s))}{c\,s^{0}}\,ds
      \le C\delta
      \qquad (n=2).
\]
Therefore for every layer budget \(\lambda\in(0,\pi/4)\) there is a regular
level \(E_t\) with
\[
   |\Om\setminus E_t|\le \lambda,
   \qquad
   D_H(t)+D_I(t)\le C\frac{\delta}{\lambda}.
   \tag{2.1}
\]
The Talenti rearrangement bound also gives
\[
   t\le C\lambda.
   \tag{2.2}
\]
Indeed, the ball level with removed area \(\lambda\) has height
\(\lambda/(4\pi)+O(\lambda^2)\), and \(u_\Om^*\le u_B^*\).

Taking
\[
   \lambda=\kappa\sqrt\delta
\]
gives a regular level \(E=\Et\) satisfying
\[
   |\Om\setminus E|\le C\sqrt\delta,\qquad
   D_H(E)+D_I(E)\le C\sqrt\delta,\qquad
   \dT(E)\le C\delta.
   \tag{2.3}
\]
The layer transfer back to \(\Om\) is still sharp:
\[
   \A(\Om)\le \A(E)+C\sqrt\delta,
   \qquad
   \A(\Om)^2\le 2\A(E)^2+C\delta.
   \tag{2.4}
\]
Thus it is enough to prove \(\A(E)^2\le C\delta\) for the
\(\sqrt\delta\)-stopped level.

\begin{remark}
This is a genuine improvement over the previous \(O(\delta)\)-layer graph
entry target.  At an \(O(\delta)\) layer, Chebyshev only gives a fixed small
\(D_H+D_I\).  At a \(\sqrt\delta\) layer, the geometric defects tend to zero.
The final exponent is not harmed because asymmetry, not volume, is squared at
the end.
\end{remark}

\section{What \(D_I\) already proves in the plane}

Write
\[
   \varepsilon:=C\sqrt\delta
\]
for the right side of \(D_I(E)\le\varepsilon\).  Since
\[
   D_I(E)=P(E)^2-4\pi |E|,
\]
and \(|E|=\pi+O(\sqrt\delta)\), after harmless rescaling we may read this as
\[
   P(E)\le 2\pi+C\varepsilon.
   \tag{3.1}
\]

\begin{lemma}[Small components are negligible]\label{lem:components}
Let \(E\subset\R^2\) have \(|E|=\pi\) and \(P(E)\le 2\pi+\varepsilon\), with
\(\varepsilon\le \varepsilon_0\).  Let \(E_0\) be an indecomposable component
of largest area and set \(E_{\rm sat}=E\setminus E_0\).  Then
\[
   |E_{\rm sat}|\le C\varepsilon^2.
\]
\end{lemma}

\begin{proof}
Let \(a_j\) be the areas of the indecomposable components, ordered so that
\(a_0\) is largest, and let \(v=\sum_{j\ge1}a_j\).  By the planar
isoperimetric inequality and additivity of perimeter over indecomposable
components,
\[
   2\pi+\varepsilon\ge P(E)
      \ge 2\sqrt\pi\left(\sqrt{\pi-v}+\sum_{j\ge1}\sqrt{a_j}\right).
\]
For small \(v\),
\[
   2\sqrt\pi\sqrt{\pi-v}
      \ge 2\pi-Cv.
\]
Hence
\[
   \sum_{j\ge1}\sqrt{a_j}\le C\varepsilon+Cv.
\]
For \(\varepsilon_0\) small, \(v\le \pi\) and the \(Cv\) term is absorbed
using \(v\le(\sum_{j\ge1}\sqrt{a_j})^2\) unless \(v\) is already bounded by
\(C\varepsilon^2\).  Thus \(v\le C\varepsilon^2\).
\end{proof}

With \(\varepsilon\simeq\sqrt\delta\), disconnected satellites have total area
\(O(\delta)\).  They can be deleted without affecting the final asymmetry
bound.  Their torsional rigidity is also negligible:
\[
   T(E_{\rm sat})\le \sum_j \frac{|E_j|^2}{4\pi}
      \le \frac{|E_{\rm sat}|^2}{4\pi}
      =O(\delta^2),
\]
by Saint--Venant on each component.

\begin{lemma}[Small holes have quadratic area]\label{lem:holes}
Let \(F\) be the main component and suppose the total perimeter of bounded
hole boundaries is at most \(p_h\).  Then the total area of those holes is
\[
   |H_{\rm holes}|\le \frac{p_h^2}{4\pi}.
\]
\end{lemma}

\begin{proof}
Apply the planar isoperimetric inequality to each hole \(H_j\):
\(|H_j|\le P(H_j)^2/(4\pi)\).  Summing and using
\(\sum_j P(H_j)^2\le(\sum_jP(H_j))^2\) gives the claim.
\end{proof}

Thus closed holes whose total boundary cost is \(O(\sqrt\delta)\) have total
area \(O(\delta)\).  This is already enough for tiny interior holes such as a
punctured disk.  The capacity mechanism from
\texttt{PLANAR\_RADIALIZATION\_PROOF\_ATTEMPT.tex} is stronger, but for the
\(\sqrt\delta\)-stopped level the perimeter estimate already gives the correct
area scale for closed holes.

\section{Convex hull and the graph envelope}

Let \(F\) be the main connected component from Lemma~\ref{lem:components}, and
let \(K=\operatorname{co}(F)\).  For planar continua,
\[
   P(K)\le P(F).
   \tag{4.1}
\]
Consequently \(K\) is a convex set with perimeter \(2\pi+O(\varepsilon)\) and
area \(\pi+O(\varepsilon)\).  By Bonnesen's inequality, after translating to a
suitable center \(z\), its inradius and circumradius satisfy
\[
   R_K-r_K\le C\sqrt\varepsilon.
   \tag{4.2}
\]
For \(\delta\) small this puts \(K\), after area normalization, in the
small-amplitude radial-graph class:
\[
   K^\sharp=\{z+r e^{i\theta}:0<r<1+\phi_K(\theta)\},
   \qquad
   \|\phi_K\|_{L^\infty}\le C\sqrt\varepsilon.
   \tag{4.3}
\]

This is a useful reduction but not yet a proof.  The convex hull can overfill
sparse outward structures.  If \(K\setminus F\) has volume merely
\(O(\varepsilon)=O(\sqrt\delta)\), the crude fill-rescale lemma gives only
\[
   \dT(K^\sharp)\le C(\dT(F)+\sqrt\delta),
\]
which is too weak.  To close the sharp result one needs the stronger statement
\[
   \dT(K^\sharp)\le C\dT(E).
   \tag{4.4}
\]
This is exactly where the remaining graph-reduction work sits.

\section{Why \(D_H\) is the right extra input}

The obstruction to using the convex hull alone is a sparse outward finger.
In the model case, take a tube of width \(w\) and length \(L\) attached to the
disk.  The convex hull may add area comparable to \(L\) times the angular span,
while the actual tube volume is only \(wL\).  Perimeter sees only \(L\), so
choosing the stopped level with \(D_I=O(\sqrt\delta)\) permits fingers of
length \(L=O(\sqrt\delta)\), but not longer.  This already makes the hull a
small-amplitude graph.

The sharp deficit comparison, however, uses the \(D_H\) weight.  On active
levels inside a planar tube,
\[
   D_H \simeq \frac{L}{w}.
   \tag{5.1}
\]
Thus a long thin tube cannot occur on a \(D_H\)-good level.  Equivalently, if
the tube survives the level \(u=t\), then either:
\begin{enumerate}[label=\arabic*.]
\item its cross-width satisfies \(w^2\gtrsim t\), since the torsion height in
a width-\(w\) tube is \(O(w^2)\); or
\item it lies below the stopped level and has already been removed.
\end{enumerate}
At the \(\sqrt\delta\)-stopped level, \(t\le C\sqrt\delta\).  Therefore any
surviving tube has \(w\gtrsim \delta^{1/4}\).  Since the torsion deficit of a
planar outward tube is \(O(wL)\), the length of a surviving tube paid by
\(\delta\) is at most
\[
   L\lesssim \frac{\delta}{w}\lesssim \delta^{3/4}.
\]
So the dangerous long-thin convexification example is not compatible with the
good-level information.

This model computation is the reason the reduction should be true in
dimension two.  It also shows why the proof must use both pieces:
\[
   D_I \quad\hbox{for perimeter/topological cleanup},\qquad
   D_H \quad\hbox{for capillary thin-finger cleanup}.
\]

\section{A sufficient missing lemma}

The following lemma would close the planar graph reduction without using BDV.

\begin{conjecture}[Good-level convexification lemma]\label{conj:convexify}
Let \(E=\Et\) be a regular \(\sqrt\delta\)-stopped good level in dimension
two, satisfying
\[
   |\Om\setminus E|\le C\sqrt\delta,\qquad
   D_H(E)+D_I(E)\le C\sqrt\delta,\qquad
   \dT(E)\le C\delta.
\]
Let \(F\) be its largest indecomposable component and \(K=\operatorname{co}(F)\).
After deleting the satellite components of total area \(O(\delta)\), there is
a star-shaped set \(G\) with
\[
   F\subset G\subset K,\qquad
   G^\sharp=\{r e^{i\theta}:0<r<1+\phi(\theta)\},
   \qquad
   \|\phi\|_{L^\infty}\le \eta_0,
\]
such that
\[
   |G\Delta E|\le C\sqrt\delta,
   \qquad
   \dT(G^\sharp)\le C\delta.
   \tag{6.1}
\]
\end{conjecture}

Assuming Conjecture~\ref{conj:convexify}, the proof closes:
\[
   \A(G)^2\le C\dT(G^\sharp)\le C\delta
\]
by the direct radial-graph theorem, while
\[
   \A(E)\le \A(G)+C|E\Delta G|\le C\sqrt\delta.
\]
Then \(\A(\Om)\le \A(E)+C\sqrt\delta\), hence \(\A(\Om)^2\le C\delta\).

\section{Attempted proof of the missing lemma}

The natural proof is a decomposition of \(K\setminus F\).

\subsection{Closed holes}

Closed holes are controlled by Lemma~\ref{lem:holes}; their area is
\(O(\delta)\).  Filling them changes asymmetry by \(O(\delta)\).  The torsion
change is harmless either by the crude fill-rescale lemma or by the stronger
capacity identity:
\[
   T(D)-T(D\setminus H)\ge a^2\operatorname{cap}_D(H)
\]
whenever \(u_D\ge a\) on \(H\).  Tiny central holes are therefore not a real
obstruction.

\subsection{Boundary fjords}

A fjord of radial depth \(h\) connected to the exterior has two sides of total
length at least \(2h\).  Thus \(D_I\le C\sqrt\delta\) excludes \(h\gg
\sqrt\delta\).  Fjords of depth \(O(\sqrt\delta)\) may be filled or ignored at
the graph-amplitude level.  If their area is \(O(\sqrt\delta)\), they sit in a
torsion boundary layer of height \(O(\sqrt\delta)\), so their torsion effect
should be \(O(\delta)\).  Making the last sentence rigorous requires a
uniform boundary-layer comparison for the convex envelope \(G\); this is not
written in the repo.

\subsection{Outward fingers}

For tube-like fingers the model estimate above proves the needed scale:
surviving fingers on the \(\sqrt\delta\)-stopped level have length
\(O(\delta^{3/4})\), hence lie inside the graph amplitude threshold, and the
ones not surviving have already been removed by the level flow.  What remains
missing is an arbitrary-geometry version:

\[
   \hbox{every non-graph outward component}
      \quad\Longrightarrow\quad
   \hbox{a family of tube estimates like }D_H\simeq L/w.
   \tag{7.1}
\]

This is not a cosmetic point.  Without (7.1), the convex hull could overfill a
sparse set of outward arcs, and the crude volume estimate would lose the sharp
rate.

\subsection{Bad rays}

Plan 2's radial slicing controls multiplicity on bad rays once an annular
bracket is known.  At the \(\sqrt\delta\)-stopped level the annular bracket is
available from the convex-hull/Bonnesen argument after deleting satellites.
The remaining problem is not measuring the bad-ray set; it is proving that the
radial repair along those bad rays changes torsion by \(O(\delta)\), not merely
that it changes area by \(O(\sqrt\delta)\).

\section{Verdict}

I cannot honestly claim the graph-reduction step is proved.  But the
\(\sqrt\delta\)-stopping scale makes the remaining theorem substantially more
plausible and more precise.

What is now proved or reduced to standard planar facts:
\begin{enumerate}[label=\arabic*.]
\item a good level with
\[
   |\Om\setminus E|\le C\sqrt\delta,\quad
   D_H(E)+D_I(E)\le C\sqrt\delta,\quad
   \dT(E)\le C\delta;
\]
\item satellite components have total area \(O(\delta)\);
\item closed holes with total perimeter \(O(\sqrt\delta)\) have total area
\(O(\delta)\);
\item the main connected component has a convex graph envelope with small
\(L^\infty\) amplitude by planar convex-hull perimeter and Bonnesen;
\item model outward fingers are removed or made harmless at the correct scale
by combining the level height \(t\le C\sqrt\delta\) with the planar tube
deficit and the \(D_H\simeq L/w\) estimate.
\end{enumerate}

The exact missing point is Conjecture~\ref{conj:convexify}: an
arbitrary-geometry decomposition showing that convexifying/radializing only
the non-graph pieces of the good level changes the normalized torsion deficit
by \(O(\delta)\).  This is narrower than the previous missing theorem and is
where I would focus next.

\end{document}
